\section{Experimental scope and inference}
\label{sec:experiments}

\subsection{Registered VOC development}
% EXP-034--040; V6 protocol, addendum, final report and handoff.
PASCAL Visual Object Classes (VOC) 2012 supplies the development data
\citep{pascal-voc-2012}. The largest
FIT panel contains \nval{voc.images} image groups and
\nval{voc.states} designed image--class states. FIT denotes this
development population; held-out fold predictions are termed out-of-fold
(OOF). Five outer folds keep every
state of an image in the same partition; scaling and early stopping use
only the corresponding training-side groups. The registered seeds are 13,
37 and 71. Continuous seed predictions are combined before selecting
actions. Seeds, folds, classes and pixels are not treated as additional
independent image groups.

The prospective study contains \nval{voc.runs} validated fits across its
capacity/representation batch, decision-objective batch and matched
F0/F1 scale comparisons. This development phase is denoted V6. The
registered cost grid is $\lambda\in\{0,0.01,0.025,0.05,0.1,0.2\}$;
the present core result tables show the zero-cost condition and preserve
the cost/budget scope in the supplement. Development sparse budgets include
0.5\%, 1\%, 2\%, 5\%, 10\% and 20\%, with 1/2/5\% designated core
budgets. Table~\ref{tab:models} records capacities and output types.
The complete-oracle zero-cost population has \nval{voc.complete} states.
The saved eight-row F0/F1 summary does not supply its exact evaluable-state
or defined-pair counts; these denominators remain unavailable in the local
source excerpt and cannot be inferred from nominal scale. The target-rich
populations are given separately in Appendix~\ref{sec:sensitivity}.

VOC FIT labels had already supported development. Registering the prospective
refit constrained its analysis but did not create a pristine held-out
dataset. We retain its frozen \texttt{MIXED} terminal and distinguish
its planned contrasts from its explicitly post-hoc composition and
alternative-utility analyses.

\subsection{Frozen COCO external transfer}
% EXP-048,050--054,056; final protocol and target/trajectory/prediction locks.
The external panel uses Common Objects in Context (COCO) val2017
\citep{lin2014coco}. A label-blind
image inventory, exact and registered strong-near duplicate checks against
the VOC FIT pool, and deterministic hash ordering fixed
\nval{coco.images} images before performance access. Every selected image
is crossed with the same vocabulary of 20 VOC classes, giving
\nval{coco.states} designed states. Class presence does not filter the
panel. VOC prompt strings remain fixed; the target construction maps them
to the registered COCO categories and unions valid instances, including
valid crowd instances. No label-derived crop, box or point is an input to
the predictor or repair plan.

Six families transfer: \model{G-small}, \model{G-match},
\model{Atom}, \model{Rect}, \model{Union} and
\model{Relative}. Their three-seed VOC full fits, scaler/schema identities,
trajectories and continuous predictions were locked before the single
recorded COCO GT read. The external study contains no COCO training or
post-target model selection. It does not externally evaluate the VOC
direct-decision baselines, F0/F1, or alternative target losses.
The executable E4 operationalization and its FORCED\_K binding were added
after prediction lock and before first GT access. Appendix~\ref{sec:freeze-timeline}
records this sequence; the evidence establishes repository-recorded
freezing, not public preregistration.

\subsection{Eligibility and technical missingness}
\label{sec:missingness-timing}
The frozen COCO accounting separates genuine empty/no-result masks from
technical execution failures. One trajectory-level failure and three A4
action-level failures affect \nval{coco.missing} states. Available,
defined action cells remain eligible for fidelity: MAE uses
\nval{coco.cells} cells and DRRE uses \nval{coco.pairs} STOP-relative
pairs across \nval{coco.mae_states} states. Complete-oracle intervention
metrics use \nval{coco.complete} states because missing a feasible outcome
prevents a complete oracle comparison. All \nval{coco.images} image groups
remain in the resampling universe. Missing values are not replaced by zero,
and failed A4 actions do not become scientifically infeasible repairs.
The saved choices on the three A4-failure states use an offline availability
mask formed after those failures, although before GT access. This handling
is label-blind but not pre-action observable. These states remain in the
defined-prediction diagnostic population; only saved non-STOP choices
contribute to its count. These states contribute neither to
complete-oracle utility nor to its sparse-budget population. They support
no deployment-utility claim for failure-aware selection
(Appendix~\ref{sec:implementation-details}).

\subsection{Paired uncertainty and frozen interpretation}
Both prospective stages use paired image-group bootstrap intervals with
\nval{coco.bootstrap} replicates, seed 13, NumPy PCG64 sampling, and linear
percentile endpoints at 0.025 and 0.975. Each resample retains the class
states and action records of a sampled image together, including
multiplicity. Complete metrics are recomputed within the original frozen
evaluation for each resample, then contrasted. The resulting intervals
support pointwise comparisons, not
multiplicity-adjusted confirmatory $p$-values.

For lower-is-better fidelity metrics, effects are left minus right.
A practical tie requires the complete paired interval to lie within
$\pm 0.05$ times the comparator's full-sample metric, with the frozen
zero-reference convention. A nonzero difference, or an interval that
excludes zero, can therefore be a practical tie. V6 additionally uses its
registered magnitude and image-group win/loss conditions for clear
improvement or worsening. COCO E1/E2 and the E4/E5 synthesis use their own
pre-target contracts; we do not exchange those stage-specific rules.
